% ===================================================================
%  TABULKA č. 8 — Žně. — Lov, vinobraní, rybolov.
%  Traduction 2026 d'apres texto/io, controlee sur texto/fr, texto/ru
%  et texto/uk. Consigne : voir texto/cs/10-tabulka-01.tex.
%
%  L'ECART DECLARE DU BLOC c2-03-1, LE DEUXIEME DES TROIS. L'ido y
%  oublie le renvoi « (13) » du gobelet, que le francais donne ; la
%  colonne le rend, avec les treize qui la precedent, et renvoji.py
%  le passe sans rien dire.
%
%  LA VENDANGE DONNE ICI CE QU'ELLE AVAIT DONNE EN GALICIEN, ET POUR
%  LA MEME RAISON. La ou l'on vendange depuis toujours, on distingue :
%  « lis » le pressoir, « matoliny » le marc, « hrozen » la grappe,
%  « révový keř » le cep, « duha » la douve, « obruč » le cercle, et
%  « sud », « káď », « kádička », « soudek » pour ce que le francais
%  nomme d'un seul « tonneau ». La Boheme et la Galice n'ont pas de
%  vignoble commun, et elles distinguent pareil.
%
%  ET LA MOISSON EN DIT AUTANT : « klas » l'epi, « snop » la gerbe,
%  « stéblo » le chaume, « brus » la pierre a aiguiser, « cep » le
%  fleau, « stoh » la meule. Six mots pour six objets qu'une seule
%  saison rassemble.
% ===================================================================

%%K t08-noto-1 noto
\VUnotes{143.2mm}{%
\textit{(*) Protože to slovo není přesné: jde o sklizeň lnu, o vinobraní, o sklizeň jablek?
Snad by bylo dobré užít slova «} \VUgras{mesono} \textit{», navrženého v Mondo XIV, s. 242.
Ale dosud ten nový kořen nebyl Akademií přijat a čeká na rozpravy podle obvyklého idistického
řádu.}%
}

%%K t08-apar-1 apar
\VUsaut{5.0mm}
\VUtitre{15.0pt}{60}{TABULKA č. 8}
\VUsaut{3.0mm}
\VUcentre{13.2pt}{90}{\textit{První scéna.}}
\VUsaut{3.0mm}
\VUpk{10.2pt}{40}{Žně (*).}
\VUsaut{4.0mm}
\VUpk{10.2pt}{40}{Podoby venkova.}
\VUsaut{3.0mm}

%%K t08-c1-01-1 p
1. --- S \VUgras{létem} se ještě jednou změnila podoba venkova. Semeno svěřené brázdě
vzklíčilo; zelená rostlinka vyšla ze země a vydala klas. Zrno se utvořilo, potom dozrálo a od
nynějška je třeba jen sklízet.

%%K t08-c1-02-1 p
2. --- \VUgras{Polní květiny} jsou rozvité. \VUgras{Chrpy} \textsuperscript{(1)},
\VUgras{vlčí máky} \textsuperscript{(2)} a \VUgras{náprstníky}
\textsuperscript{(3)} mísí do jednotvárného odstínu obilných polí veselý kontrast svých
svěžích \VUgras{korun.}

%%K t08-c1-03-1 p
3. --- Ovocné stromy se pokryly listím; ovoce dozrálo. Malá \VUgras{třešeň}
\textsuperscript{(4)} oplývá krásnými \VUgras{třešněmi} \textsuperscript{(5)}, které děti s
velkou rozkoší trhají a jedí.

%%K t08-c1-04-1 p
4. --- Od nynějška může dobytek trávit celý den na volném vzduchu.

%%K t08-c1-04-2 p
\VUgras{Jehňata} \textsuperscript{(6)} vyrostla a stala se z nich krásné \VUgras{ovce}
\textsuperscript{(7)} a krásní \VUgras{skopci} \textsuperscript{(8)} s hustým vlněným rounem.
Klidně spásají \VUgras{trávu} \VUgras{pastviny} \textsuperscript{(9)} pestré
\VUgras{sedmikráskami.}

%%K t08-c1-04-3 p
Brzy se \VUgras{ta} zvířata začnou stříhat, aby se získala jejich \VUgras{vlna,} z níž se
vyrábí sukno našich šatů.

%%K t08-tit-2 sub
\VUsaut{5.0mm}
\VUpk{10.2pt}{40}{Pastýř.}
\VUsaut{3.4mm}

%%K t08-c1-05-1 p
5. --- \VUgras{Pastýř} \textsuperscript{(10)} si jich, zdá se, příliš nevšímá. Stará se jen o
bouři, která přechází na obzoru, a \VUgras{ovčácký hlídač} \textsuperscript{(13)}, který si
přikládá ke rtům \VUgras{polní láhev} \textsuperscript{(14)}, se zdá ještě bezstarostnější.

%%K t08-c1-06-1 p
6. --- Vedro je tísnivé, protože i \VUgras{pes} \textsuperscript{(15)} přichází uhasit žízeň;
chlemtá svěží vodu \VUgras{potoka} \textsuperscript{(16)} a plaší ubohou
\VUgras{žabku} \textsuperscript{(17)}, která se schoulila v pobřežní trávě. Ta skáče do
čiré vody, kde kvetou \VUgras{lekníny} \textsuperscript{(18)}.

%%K t08-c1-07-1 p
7. --- \VUgras{Konipas} \textsuperscript{(19)} se koupe u \VUgras{pramínku}
\textsuperscript{(20)}, který tryská mezi dvěma balvany, a ukrývá se pod
\VUgras{trnitými keři} \textsuperscript{(21)} a \VUgras{hlohovými houštinami}
\textsuperscript{(22)}.

%%K t08-tit-3 sub
\VUsaut{5.0mm}
\VUpk{10.2pt}{40}{Žně.}
\VUsaut{3.4mm}

%%K t08-c1-08-1 p
8. --- Pro venkovské děti jsou žně velmi zábavným obdobím. Vesele se
\VUgras{kotrmelcují} \textsuperscript{(24)} na \VUgras{hromadách slámy}
\textsuperscript{(23)}, pomáhají \VUgras{žencům} \textsuperscript{(26)} a, zatímco ti sekají
\VUgras{obilné pole} \textsuperscript{(27)} svými \VUgras{kosami} \textsuperscript{(25)} dobře
nabroušenými \VUgras{brusem} \textsuperscript{(30)}, sbírají obilí a vážou je do
\VUgras{snopů} \textsuperscript{(31)} nebo do
\VUgras{otepí} \textsuperscript{(32)}.

%%K t08-c1-09-1 p
9. --- Někdy se těm největším z nich dovolí řezat \VUgras{srpem} \textsuperscript{(12)}
\VUgras{zlatá stébla} \textsuperscript{(28)}, která nesou
\VUgras{klasy} \textsuperscript{(29)} plné zrna; a ti malí, hlídajíce \VUgras{dobytek}
\textsuperscript{(33)}, který se pase na \VUgras{pastvině} \textsuperscript{(34)} ve stínu
velké \VUgras{lípy} \textsuperscript{(35)}, chytají svou \VUgras{síťkou}
\textsuperscript{(36)} krásné \VUgras{motýly.}

%%K t08-c1-09-2 p
Někteří dokonce lezou na \VUgras{vozík} \textsuperscript{(37)}, aby urovnali snopy, které jim
podávají \VUgras{vidlemi} \textsuperscript{(38)}.

%%K t08-c1-10-1 p
10. --- Na celé \VUgras{rovině} \textsuperscript{(39)}, kde vzlétají \VUgras{skřivani}
\textsuperscript{(41)}, vládne ruch. Všichni jsou zaměstnáni žněmi. Ale zatímco chudí dále
užívají starých nástrojů, \VUgras{kos, srpů} a \VUgras{cepů} \textsuperscript{(40)}, bohatí
statkáři dávají sekat své obilí nebo své \VUgras{žito} \textsuperscript{(43)}
\VUgras{žacím strojem} \textsuperscript{(42)}. Potom, protože není třeba odvážet snopy do
stodoly, dělají se velké \VUgras{stohy} \textsuperscript{(44)}.

%%K t08-c1-11-1 p
11. --- Tehdy se přiváží \VUgras{mlátička} \textsuperscript{(47)}, kterou pohání
\VUgras{lokomobila} \textsuperscript{(45)} \VUgras{hnacím řemenem}
\textsuperscript{(46)}, a tak se zrno odděluje od \VUgras{slámy,} aniž opustí pole, kde bylo
zaseto.

%%K t08-tit-4 sub
\VUsaut{5.0mm}
\VUpk{10.2pt}{40}{Bouře.}
\VUsaut{3.4mm}

%%K t08-c1-12-1 p
12. --- Často se v době žní strhnou \VUgras{bouřky} \textsuperscript{(53)}. Nejprve je z dálky
slyšet dunění \VUgras{hromu}; \VUgras{západní vichr} \textsuperscript{(54)} začíná vát a
udýchaně ohýbá nejtlustší stromy \VUgras{lesa} \textsuperscript{(49)}. Černé \VUgras{mraky}
\textsuperscript{(56)} útočí na nebe; prostupují je oslnivé klikatiny \VUgras{blesku}
\textsuperscript{(55)}. \VUgras{Déšť} a někdy \VUgras{krupobití} začínají padat a drtí úrodu
připravenou k posekání.

%%K t08-c1-13-1 p
13. --- Často blesk zapálí stavení. Tak loni byla střecha \VUgras{větrného mlýna}
\textsuperscript{(50)} obrácena v popel a dvě jeho \VUgras{křídla} \textsuperscript{(51)}
rozbita.

%%K t08-c1-14-1 p
14. --- Po několika hodinách se znovu rodí klid. Zářivá \VUgras{duha} \textsuperscript{(52)}
se objevuje na obzoru a příroda se ujímá svého života přerušeného na několik okamžiků.
Venkované, kteří se uchýlili do \VUgras{chatrčí} \textsuperscript{(48)}, se vracejí k práci a
odvážně se snaží napravit škody, které právě utrpěli.

%%K t08-tit-5 sub
\VUsaut{6.0mm}
\VUcentre{13.2pt}{90}{\textit{Druhá scéna.}}
\VUsaut{3.6mm}

%%K t08-tit-6 sub
\VUpk{10.2pt}{40}{Lov. --- Vinobraní.}
\VUsaut{1.4mm}
\VUpk{10.2pt}{40}{Rybolov.}
\VUsaut{4.0mm}

%%K t08-c2-01-1 p
1. --- Ke konci \VUgras{srpna} nebo v polovině \VUgras{září}, podle krajů, se zahajuje lov.
Sotva první sluneční paprsky vyhnaly \VUgras{ještěrky} \textsuperscript{(2)} z jejich díry, už
se \VUgras{lovec} \textsuperscript{(5)} vydává na cestu se svými
\VUgras{psy} \textsuperscript{(4)}.

%%K t08-c2-02-1 p
2. --- Oblékl si svůj pevný \VUgras{lovecký oblek} \textsuperscript{(9)} z tlustého sukna,
obul si kožené \VUgras{kamaše,} s nimiž se mu podaří chodit ve vlhké trávě, kde se skrývají
\VUgras{ropuchy} \textsuperscript{(1)}; vzal svou \VUgras{pušku} \textsuperscript{(6)} a svou
\VUgras{brašnu na zvěř} \textsuperscript{(10)}, a už odchází.

%%K t08-c2-03-1 p
3. --- Horlivost pronásledování ho přivádí doprostřed vinice, kde je vinobraní velmi čilé.
Vinařovy děti, které obyčejně hlídají kozy na cestě, je dnes nechávají pást nazdařbůh a,
přivázavše ke kůlu starého \VUgras{kozla} \textsuperscript{(3)}, který by neopomněl využít své
svobody k útěku, přiznávají si i ony podíl na potěšení. Jeden uchopí hrozen v
\VUgras{koši na vinobraní} \textsuperscript{(12)} a baví se tím, že nechá téci \VUgras{šťávu}
\textsuperscript{(14)} do \VUgras{poháru} \textsuperscript{(13)}, aby vyrobil mošt (nezkvašené
víno). Druhý si nasazuje
\VUgras{věnec z listí} \textsuperscript{(15)}, který právě upletl. Blízko nich
přilétají nestydatí \VUgras{vrabci} \textsuperscript{(16)} až před lis, aby uloupili hrozny.

%%K t08-c2-04-1 p
4. --- Zatímco se děti baví, muži pracují. \VUgras{Vinaři} \textsuperscript{(18)} přinášejí
hrozny do sklepa v \VUgras{nůších} \textsuperscript{(17)}; \VUgras{bednář}
\textsuperscript{(19)} opravuje \VUgras{sudy} \textsuperscript{(20)}, ověřuje, zda jsou
\VUgras{duhy} \textsuperscript{(22)} a
\VUgras{dna} \textsuperscript{(21)} v pořádku, a dává nahradit prasklé
\VUgras{obruče} \textsuperscript{(23)}. Pracoval také na \VUgras{kádích}
\textsuperscript{(25)}, do nichž se sypou \VUgras{hrozny} \textsuperscript{(26)}, aby kvasily.

%%K t08-c2-05-1 p
5. --- Když kvašení skončí, víno se stočí a matoliny se položí na \VUgras{lis}
\textsuperscript{(28)} a postupným utahováním velkého \VUgras{šroubu} \textsuperscript{(29)}
se vytlačí šťáva, která teče do \VUgras{kádičky} \textsuperscript{(32)}, a získá se ještě
\VUgras{víno} \textsuperscript{(31)}.

%%K t08-tit-8 sub
\VUsaut{6.0mm}
\VUpk{10.2pt}{40}{Pivo a jablečné víno.}
\VUsaut{3.4mm}

%%K t08-c2-06-1 p
6. --- Na zdi \VUgras{sklepa} \textsuperscript{(27)} šplhá větev \VUgras{chmele}
\textsuperscript{(30)}. Tam je jen okrasnou rostlinou, ale v některých zemích, kde je réva
velmi vzácná, se chmel pěstuje k výrobě \VUgras{piva.}

%%K t08-c2-07-1 p
7. --- Malá \VUgras{jabloň} \textsuperscript{(33)} je obtěžkána jablky, která dají, až budou
zralá, výborné \VUgras{jablečné víno.} Ve své \VUgras{kleci} \textsuperscript{(34)} hledí
\VUgras{kanárek} \textsuperscript{(35)} a \VUgras{stehlík} \textsuperscript{(36)} závistivýma
očima na vrabce, kteří volně poskakují.

%%K t08-c2-08-1 p
8. --- Kam až oko dohlédne, na svahu pahorků jsou v řadách postaveny \VUgras{vinné kůly}
\textsuperscript{(40)}, které podpírají překrásné \VUgras{révové keře} \textsuperscript{(39)},
obtěžkané těžkými \VUgras{hrozny.}

%%K t08-c2-08-2 p
Po celý rok \VUgras{vinař} \textsuperscript{(42)} pracoval, aby pečoval o svou
\VUgras{vinici} \textsuperscript{(38)}, a nyní je za svou námahu odměněn krásnou
sklizní. \VUgras{Vinařky} \textsuperscript{(41)} plní kádě postavené na vozíku taženém
\VUgras{mezky} \textsuperscript{(43)}, a všichni jsou veselí.

%%K t08-tit-9 sub
\VUsaut{5.0mm}
\VUpk{10.2pt}{40}{Rybolov.}
\VUsaut{3.4mm}

%%K t08-c2-09-1 p
9. --- Stejně je tomu s \VUgras{rybáři na udici} \textsuperscript{(44)}, kteří se od rána
přišli usadit na břehu vody se svými \VUgras{udicemi} \textsuperscript{(45)}.

%%K t08-c2-09-2 p
\VUgras{Ryba} \textsuperscript{(47)} chňapne po \VUgras{háčku,} jakmile ponoří svůj
\VUgras{vlasec} \textsuperscript{(46)} do vody, a sotva měli čas obnovit
\VUgras{návnadu,} už se nová oběť zmítá na konci vlasce.

%%K t08-c2-09-3 p
Přesto \VUgras{rybář se sítí} \textsuperscript{(48)}, který ze své \VUgras{loďky}
\textsuperscript{(49)} vrhá \VUgras{vrhací síť} doprostřed \VUgras{řeky}
\textsuperscript{(50)}, chytá více ryb.

%%K t08-c2-10-1 p
10. --- Podzim je obdobím pěkných vycházek: pěšky, \VUgras{na koni} \textsuperscript{(52)},
\VUgras{na kole} \textsuperscript{(53)}, \VUgras{autem} \textsuperscript{(54)}. \VUgras{Cesty}
\textsuperscript{(51)} jsou čisté a užívá se posledních krásných dní, protože už se
\VUgras{vlaštovky} \textsuperscript{(56)} shromažďují na střechách, aby ze všech sil odlétly
do teplejších zemí. Listí starého \VUgras{ořešáku} \textsuperscript{(57)} žloutne,
\VUgras{ořechy} padají a vysoké \VUgras{stromy} seskupené jako \VUgras{kytice}
\textsuperscript{(58)} na \VUgras{výšině} \textsuperscript{(59)} brzy ztratí listí.

%%K t08-c2-11-1 p
11. --- \VUgras{Slunce} \textsuperscript{(60)} vychází později, dny se krátí, dosti ostrý
chlad se začíná pociťovat \VUgras{ráno,} a už hejna \VUgras{sluk} \textsuperscript{(61)}
opouštějí naše nehostinné podnebí.
